\documentclass{umk-matematika}
\begin{document}

\begin{center}
{\Large\bfseries Практическая работа 18}
\end{center}
\vspace{0.5cm}

\textbf{Задача 1.} Решите неравенство: $2^x > 8$\par\vspace{10pt}
\textbf{Задача 2.} Решите неравенство: $3^{x-2} \le 27$\par\vspace{10pt}
\textbf{Задача 3.} Решите неравенство: $\left(\frac{1}{2}\right)^x < 4$\par\vspace{10pt}
\textbf{Задача 4.} Решите неравенство: $5^{2x} \ge 1$\par\vspace{10pt}
\textbf{Задача 5.} Решите неравенство: $4^{x} - 6 \cdot 2^{x} + 8 \le 0$\par\vspace{10pt}
\textbf{Задача 6.} Решите неравенство: $9^{x} - 10 \cdot 3^{x} + 9 > 0$\par\vspace{10pt}
\textbf{Задача 7.} Решите неравенство: $2^{x+1} + 2^{2-x} \le 9$\par\vspace{10pt}
\textbf{Задача 8.} Количество кирпичей на стройке удваивается каждый день по закону $N(t) = 100 \cdot 2^{t}$, где $t$ — дни. Через сколько дней количество кирпичей превысит 3200 штук?\par\vspace{10pt}
\textbf{Задача 9.} Давление в трубе падает по закону $P(h) = 64 \cdot \left(\frac{1}{2}\right)^{\frac{h}{2}}$, где $h$ — высота в метрах. На какой высоте давление станет меньше 4 единиц?\par\vspace{10pt}
\textbf{Задача 10.} Решите неравенство: $\frac{3^{x} - 9}{2^{x} - 4} \ge 0$\par\vspace{10pt}
\newpage
\section*{Ответы}

\begin{enumerate}
  \item Ответ: $x \in (3; +\infty)$
  \item Ответ: $x \in (-\infty; 5]$
  \item Ответ: $x \in (-2; +\infty)$
  \item Ответ: $x \in [0; +\infty)$
  \item Ответ: $x \in [1; 2]$
  \item Ответ: $x \in (-\infty; 0) \cup (2; +\infty)$
  \item Ответ: $x \in [-1; 2]$
  \item Ответ: через $6$ дней (на 6-й день количество превысит 3200)
  \item Ответ: $h > 8$ метров
  \item Ответ: $x \in (-\infty; 2) \cup (2; +\infty)$
\end{enumerate}
\end{document}
